\documentclass[12pt, a4paper]{report}

% Packages
\usepackage[utf8]{inputenc}
\usepackage[T1]{fontenc}
\usepackage[english]{babel}
\usepackage{amsmath, amssymb, amsthm, amsfonts}
\usepackage{geometry}
\usepackage{hyperref}
\usepackage{bookmark}
\usepackage{enumitem}
\usepackage{fancyhdr}
\usepackage{titlesec}
\usepackage{mathrsfs}
\usepackage{xcolor}
\usepackage{tikz}
\usepackage{pgfplots}
\pgfplotsset{compat=1.18}

% Geometry setup
\geometry{
    top=2.5cm,
    bottom=2.5cm,
    left=2.5cm,
    right=2.5cm
}
\setlength{\headheight}{30pt}

% Hyperref setup
\hypersetup{
    colorlinks=true,
    linkcolor=blue,
    filecolor=magenta,      
    urlcolor=cyan,
    pdftitle={Stochastic Finance and Risk Modelling - Chapter 7},
    pdfauthor={Laurence Carassus and Gaoyue Guo},
}

% Theorem environments
\newtheorem{theorem}{Theorem}[chapter]
\newtheorem{proposition}[theorem]{Proposition}
\newtheorem{lemma}[theorem]{Lemma}
\newtheorem{corollary}[theorem]{Corollary}

\theoremstyle{definition}
\newtheorem{definition}{Definition}[chapter]
\newtheorem{example}{Example}[chapter]
\newtheorem{remark}[theorem]{Remark}

% Layout/Style
\pagestyle{fancy}
\fancyhf{}
\rhead{\leftmark}
\lhead{}
\cfoot{\thepage}

\title{\textbf{Stochastic Finance and Risk Modelling} \\ \large Chapter 7: American Options}
\author{Based on slides by Ioane Muni-Toke \\ Laurence Carassus and Gaoyue Guo \\ \textit{Expanded Study Resource}}
\date{\today}

\begin{document}

\maketitle

\tableofcontents

\chapter*{Introduction}
\addcontentsline{toc}{chapter}{Introduction}
This document serves as a comprehensive study guide for the course "Stochastic Finance and Risk Modelling," functioning as a complete substitute for the lecture slides related to Chapter 7: "American Options."

Unlike European options, which can only be exercised at maturity $T$, American options provide the holder with the right to exercise the option at any time up to and including the maturity date. This early exercise feature adds significant complexity to the pricing and hedging of these instruments. To mathematically formalize the holder's choice of when to exercise, we rely heavily on the theory of \textit{stopping times} and the \textit{Snell envelope}, which solves the underlying optimal stopping problem.

In this chapter, we first formalize the American option setup. We then rigorously define stopping times, stopped processes, and their properties. Following this, we introduce the crucial concept of the Snell envelope as the smallest supermartingale dominating the payoff process, linking it explicitly to optimal stopping times. Finally, we establish the pricing equations for American options, compare them to their European counterparts, and discuss replication strategies from both the buyer and seller perspectives in a complete market.

Detailed explanations, rigorous mathematical derivations, and intuition are provided throughout to bridge conceptual gaps and ensure a deep understanding of the material.

\chapter{First Approach}

In this chapter, we develop the mathematical framework necessary to price an American option in discrete time. We will start with standard definitions and a first look at the recursive pricing approach, often referred to as dynamic programming.

\section{Model and Definitions}

\begin{definition}[Financial Market Framework]
We consider a financial market in discrete time $t \in \{0, 1, \dots, T\}$. 
The market consists of a riskless asset (e.g., a bank account) denoted by its price process $(S^0_t)_{0 \leq t \leq T}$ and $d=1$ risky asset with price process $(S_t)_{0 \leq t \leq T}$. 

The uncertainty in the market is modeled on a rigorous probability space $(\Omega, \mathcal{F}, \mathbb{P})$, which is endowed with a natural filtration $\mathbb{F} = (\mathcal{F}_{t})_{0 \leq t \leq T}$. 
The $\sigma$-algebra $\mathcal{F}_t$ represents all the information available in the market up to and including time $t$. 
For simplicity and mathematical standard practice, we assume that the initial state is deterministic or completely known, meaning $\mathcal{F}_{0} = \{\emptyset, \Omega\}$ (the trivial $\sigma$-algebra). We also assume that $\mathcal{F}_{T} = \mathcal{F}$, meaning all uncertainty is resolved by the final maturity $T$.
\end{definition}

Our focus is exclusively on the pricing of \textit{American options}. Recall that the defining characteristic of an American option is its early exercise capability. 

\begin{definition}[American Option Payoff]
The holder of an American option can choose to exercise it at \textit{any} time $t \in \{0, \dots, T\}$. If the option is exercised exactly at time $t$, the holder will receive a payoff equal to $H_{t}$.
Here, $H_{t}$ is assumed to be a positive random variable that is $\mathcal{F}_{t}$-measurable. This measurability condition ensures that the payoff dimensioned at time $t$ only depends on market information available up to time $t$ (there is no prescience of the future). The option can only be exercised exactly once.
\end{definition}

\begin{example}[American Put Option]
The quintessential example is the American put option of strike $K$ on the risky asset $S_t$. The intrinsic payoff if exercised at time $t$ is defined as:
$$
H_t = \max(K - S_t, 0) = (K - S_t)^+
$$
\end{example}

For the entirety of this pricing formulation, we operate under the essential pillars of modern mathematical finance.
\begin{enumerate}
    \item \textbf{No-Arbitrage:} We assume the market is free of arbitrage opportunities. By the First Fundamental Theorem of Asset Pricing, this implies the existence of an \textit{equivalent martingale measure} (EMM), denoted by $\mathbb{Q}$.
    \item \textbf{Completeness:} We assume the market is complete. By the Second Fundamental Theorem of Asset Pricing, the equivalent martingale measure $\mathbb{Q}$ is \textit{unique}. Hence, any derivative payoff can be perfectly replicated by a trading strategy.
\end{enumerate}

\section{Dynamic Programming Formulation}

To find the fair value of an American option at any time $t$, denoted by $U_t$, we introduce a backward induction approach, also known as the dynamic programming principle.

\begin{proposition}[Dynamic Pricing Equation]
Let $U_t$ be the value of the American option at time $t$. Then the value process sequence $(U_t)_{t=0, \dots, T}$ perfectly satisfies the following backward recursion:
$$
\begin{cases}
U_T = H_T \\
U_{t-1} = \max\left( H_{t-1}, \, S_{t-1}^0 \mathbb{E}^{\mathbb{Q}}\left[ \frac{U_t}{S_t^0} \Big| \mathcal{F}_{t-1} \right] \right), \quad t = 1, \dots, T
\end{cases}
$$
\end{proposition}

\textbf{Intuition and Explanation:}
This formulation is the heart of American option pricing. Let us break it down logically:
\begin{itemize}
    \item \textbf{At maturity ($t = T$):} The holder has no more time value left. They must either exercise the option to receive its intrinsic value $H_T$ or let it expire worthless. Thus, the value at maturity is strictly the payoff itself: $U_T = H_T$.
    \item \textbf{At any prior time ($t-1$):} The holder must make a rational choice between two mutually exclusive actions:
    \begin{enumerate}
        \item \textbf{Exercise immediately:} The holder receives the intrinsic payoff available right now, which is $H_{t-1}$.
        \item \textbf{Hold (continue):} If the holder decides \textit{not} to exercise, they hold onto the option for at least one more time step. The value of holding the option is the expected present value of the option's worth tomorrow ($U_t$), under the risk-neutral measure $\mathbb{Q}$, discounted back to today. This represents the \textit{continuation value}: 
        $$ C_{t-1} = S_{t-1}^0 \mathbb{E}^{\mathbb{Q}}\left[ \frac{U_t}{S_t^0} \Big| \mathcal{F}_{t-1} \right] $$
    \end{enumerate}
    Assuming rational behavior, the holder will always choose the action that maximizes their payoff. Therefore, the value of the option today, $U_{t-1}$, must be the maximum of the immediate exercise value and the continuation value.
\end{itemize}

\subsection{Supermartingale Property of the Discounted Value}

Let us define the \textit{discounted value} of the option and the payoff, a standard transformation that simplifies martingale analysis. We denote the discounted option value process as $\tilde{U}_t = U_t / S_t^0$ and the discounted payoff process as $\tilde{H}_t = H_t / S_t^0$.

By rewriting the dynamic programming equation in terms of discounted quantities, we get:
$$
\tilde{U}_{t-1} = \max\left( \tilde{H}_{t-1}, \, \mathbb{E}^{\mathbb{Q}}\left[ \tilde{U}_t \Big| \mathcal{F}_{t-1} \right] \right)
$$

This immediately reveals a critical theoretical property of the American option price process under the risk-neutral measure.

\begin{proposition}[Supermartingale Characterization]
The discounted value process $(\tilde{U}_t)_{t=0, \dots, T}$ of the American option is a $\mathbb{Q}$-supermartingale. Furthermore, it is the \textit{smallest} supermartingale that dominates the discounted payoff process $(\tilde{H}_t)_{t=0, \dots, T}$ from above.
\end{proposition}

\begin{proof}
By definition of a supermartingale, we need to show that for any time $t \in \{1,\dots,T\}$, 
$$ \tilde{U}_{t-1} \geq \mathbb{E}^{\mathbb{Q}}\left[ \tilde{U}_t \Big| \mathcal{F}_{t-1} \right] $$ 

Looking directly at the discounted dynamic programming equation:
$$
\tilde{U}_{t-1} = \max\left( \tilde{H}_{t-1}, \, \mathbb{E}^{\mathbb{Q}}\left[ \tilde{U}_t \Big| \mathcal{F}_{t-1} \right] \right)
$$
Because the maximum of two quantities must be greater than or equal to either of those individual quantities, it strictly follows that:
$$ \tilde{U}_{t-1} \geq \mathbb{E}^{\mathbb{Q}}\left[ \tilde{U}_t \Big| \mathcal{F}_{t-1} \right] $$
thus proving that $(\tilde{U}_t)$ is a $\mathbb{Q}$-supermartingale.

Also, it is evident from the same $\max()$ function that:
$$ \tilde{U}_{t-1} \geq \tilde{H}_{t-1} \quad \text{for all } t $$
Hence, the discounted value $\tilde{U}_t$ always bounds the discounted intrinsic payoff $\tilde{H}_t$ from above.

Finally, any other $\mathbb{Q}$-supermartingale $M_t$ that also dominates $\tilde{H}_t$ from above (i.e., $M_t \geq \tilde{H}_t$) will, by backward induction, be greater than or equal to $\tilde{U}_t$, proving that $\tilde{U}_t$ is the \textit{smallest} such supermartingale. We will formalize this concept rigorously in the upcoming sections under the terminology of the \textit{Snell envelope}.
\end{proof}

\chapter{Stopping Time and Stopped Processes}

To rigorously define the action of ''choosing when to exercise'' an American option, we must introduce the mathematical concept of a stopping time. Because an investor relies strictly on information obtained up to the present moment, they cannot base their exercise decision on future events (i.e., they possess no clairvoyance).

\section{Stopping Time}

\begin{definition}[Stopping Time]
A random variable $\tau : \Omega \to \{0, 1, \dots, T\} \cup \{+\infty\}$ is called an $\mathbb{F}$-stopping time if for every integer $t \in \{0, \dots, T\}$, the event that the stopping time has occurred by time $t$ belongs to the information set available at time $t$. Mathematically, this is expressed as:
$$ \{\tau \leq t\} \in \mathcal{F}_t \quad \text{for every } t \in \{0, 1, \dots, T\}. $$
\end{definition}

\textbf{Interpretation (No-anticipation):}
This elegantly formalizes the principle of no-anticipation. If an investor uses a trading strategy, they must decide whether to stop (or exercise) at time $t$ using only the information available in the market up to time $t$, encapsulated by the $\sigma$-algebra $\mathcal{F}_t$. They are forbidden from peeking at $\mathcal{F}_{t+1}$ to make a decision at time $t$.

\begin{proposition}[Properties of Stopping Times]
The following fundamental properties hold true:
\begin{enumerate}[label=(\alph*)]
    \item Every deterministic constant time $t \in \{0, \ldots, T\}$ is trivially a stopping time.
    \item A random variable $\tau$ taking values in $\{0, 1, \dots, T\}$ is a stopping time if and only if its precise occurrence event is measurable at that time:
    $$ \{\tau = t\} \in \mathcal{F}_t \quad \text{for every } t = 0, \ldots, T. $$
    \item If $\tau_1$ and $\tau_2$ are both stopping times, their minimum $\tau_1 \wedge \tau_2 = \min(\tau_1, \tau_2)$ and their maximum $\tau_1 \vee \tau_2 = \max(\tau_1, \tau_2)$ are also valid stopping times.
\end{enumerate}
\end{proposition}

\begin{proof}
Let us provide a detailed proof for these essential properties.
\begin{enumerate}[label=(\alph*)]
    \item Let $\tau(\omega) = t_0$ for all $\omega \in \Omega$, where $t_0$ is a constant. Then, for any $t$:
    $$ 
    \{\tau \leq t\} = 
    \begin{cases} 
    \Omega & \text{if } t_0 \leq t \\
    \emptyset & \text{if } t_0 > t
    \end{cases}
    $$
    Since both $\Omega$ and $\emptyset$ belong to every $\sigma$-algebra $\mathcal{F}_t$, the condition $\{\tau \leq t\} \in \mathcal{F}_t$ is unconditionally met.

    \item Let us prove the equivalence (\textbf{if and only if}):
    \begin{itemize}
        \item $(\implies)$ Assume $\tau$ is a stopping time, so $\{\tau \leq t\} \in \mathcal{F}_t$. Then, the event $\{\tau = t\} = \{\tau \leq t\} \setminus \{\tau \leq t-1\}$.
        Since $\mathbb{F}$ is a filtration, $\mathcal{F}_{t-1} \subseteq \mathcal{F}_t$, meaning $\{\tau \leq t-1\} \in \mathcal{F}_{t-1} \subseteq \mathcal{F}_t$. Since $\sigma$-algebras are closed under set difference, $\{\tau \leq t\} \setminus \{\tau \leq t-1\} \in \mathcal{F}_t$. Hence, $\{\tau = t\} \in \mathcal{F}_t$.
        \item $(\impliedby)$ Assume $\{\tau = k\} \in \mathcal{F}_k$ for all $k$. We want to show $\{\tau \leq t\} \in \mathcal{F}_t$. Note that:
        $$ \{\tau \leq t\} = \bigcup_{k=0}^t \{\tau = k\} $$
        For every $k \leq t$, we have $\{\tau = k\} \in \mathcal{F}_k \subseteq \mathcal{F}_t$. Since $\sigma$-algebras are closed under countable unions, it follows that their union must belong to $\mathcal{F}_t$. Thus, $\{\tau \leq t\} \in \mathcal{F}_t$.
    \end{itemize}

    \item Consider the minimum $\tau_1 \wedge \tau_2$. For $\{ \tau_1 \wedge \tau_2 > t \}$, this occurs if and only if both stopping times are strictly greater than $t$:
    $$ \{\tau_1 \wedge \tau_2 > t\} = \{\tau_1 > t\} \cap \{\tau_2 > t\} $$
    Taking the complement, we have:
    $$ \{\tau_1 \wedge \tau_2 \leq t\} = \{\tau_1 \leq t\} \cup \{\tau_2 \leq t\} $$
    Since both $\tau_1$ and $\tau_2$ are stopping times, both sets on the RHS are in $\mathcal{F}_t$. The union is also in $\mathcal{F}_t$. Thus $\tau_1 \wedge \tau_2$ is a stopping time.

    For the maximum $\tau_1 \vee \tau_2$, this is even more straightforward:
    $$ \{\tau_1 \vee \tau_2 \leq t\} = \{\tau_1 \leq t\} \cap \{\tau_2 \leq t\} $$
    The intersection of two sets in $\mathcal{F}_t$ is also in $\mathcal{F}_t$. Thus $\tau_1 \vee \tau_2$ is a stopping time.
\end{enumerate}
\end{proof}

\section{Stopped Processes (Processus Arrêté)}

If we monitor a stochastic process $(X_t)_{0 \leq t \leq T}$ (like an option value or a stock price) and we decide to halt our observation or trading at a specific random stopping time $\tau$, we create a new mathematically defined process called the stopped process.

\begin{definition}[Stopped Process]
Let $\tau$ be a stopping time and let $X = (X_t)_{t \in \{0, \dots, T\}}$ be a stochastic process. The process stopped at $\tau$, denoted by $X^\tau = (X_{t \wedge \tau})_{0 \leq t \leq T}$, is defined point-wise for all $\omega \in \Omega$ as:
$$
X_{t \wedge \tau}(\omega) =
\begin{cases}
X_t(\omega) & \text{if } t < \tau(\omega) \\
X_{\tau(\omega)}(\omega) & \text{if } t \geq \tau(\omega)
\end{cases}
$$
\end{definition}
This essentially means that the process evolves normally up to the random time $\tau$, after which it is ''frozen'' and takes a constant value equal to the value the process had at precisely the stopping time instance $X_\tau$.

One of the most powerful theorems relating to stopping times revolves around the preservation of martingale properties. Doob's Optional Stopping Theorem (for bounded variables) assures us that martingales have no deterministic drift over time, even when stopped randomly, as long as the stopping rule does not ''cheat'' by looking into the future.

\begin{proposition}[Preservation of Martingale Properties]
Let $(X_t)_{t=0, \dots, T}$ be a stochastic process and let $\tau$ be a stopping time. Then the stopped process $(X_{t \wedge \tau})_{t=0, \dots, T}$ inherits the martingale properties of the original process. Specifically:
\begin{enumerate}[label=(\alph*)]
    \item If $(X_t)_t$ is an $\mathbb{F}$-adapted process, then the stopped process $(X_{t \wedge \tau})_t$ remains $\mathbb{F}$-adapted.
    \item If $(X_t)_t$ is a martingale, the stopped process $(X_{t \wedge \tau})_t$ is a martingale.
    \item If $(X_t)_t$ is a submartingale, the stopped process $(X_{t \wedge \tau})_t$ is a submartingale.
    \item If $(X_t)_t$ is a supermartingale, the stopped process $(X_{t \wedge \tau})_t$ is a supermartingale.
\end{enumerate}
\end{proposition}

\begin{proof}
While a fully rigorous general proof involves indicator functions and optional sampling paradigms, we can construct an intuitive algebraic proof for discrete time.

First, observe that we can rewrite the increment of the stopped process algebraically:
$$ X_{t \wedge \tau} - X_{(t-1) \wedge \tau} = (X_t - X_{t-1}) \mathbf{1}_{\{\tau \geq t\}} $$
Why? If $\tau \geq t$, the stopped process has not frozen yet, so both side match: $X_t - X_{t-1}$. 
If $\tau < t$ (i.e. $\tau \leq t-1$), the process has already frozen. Thus $X_{t \wedge \tau} = X_{\tau}$ and $X_{(t-1) \wedge \tau} = X_\tau$, making the increment $0$, matching the right side where the indicator is $0$.

Now, let's take the conditional expectation of the stopped process:
\begin{align*}
\mathbb{E}\left[ X_{t \wedge \tau} \mid \mathcal{F}_{t-1} \right] &= \mathbb{E}\left[ X_{(t-1) \wedge \tau} + (X_t - X_{t-1}) \mathbf{1}_{\{\tau \geq t\}} \mid \mathcal{F}_{t-1} \right] \\
&= X_{(t-1) \wedge \tau} + \mathbb{E}\left[ (X_t - X_{t-1}) \mathbf{1}_{\{\tau \geq t\}} \mid \mathcal{F}_{t-1} \right] \\
\end{align*}
Because $\tau$ is a stopping time, the event $\{\tau \geq t\}$ is exactly the complement of $\{\tau \leq t-1\}$. Since $\{\tau \leq t-1\} \in \mathcal{F}_{t-1}$, its complement $\{\tau \geq t\}$ is also in $\mathcal{F}_{t-1}$. 
Therefore, the indicator function $\mathbf{1}_{\{\tau \geq t\}}$ is $\mathcal{F}_{t-1}$-measurable. By standard properties of conditional expectation, it can be pulled out of the expectation:
$$
\mathbb{E}\left[ X_{t \wedge \tau} \mid \mathcal{F}_{t-1} \right] = X_{(t-1) \wedge \tau} + \mathbf{1}_{\{\tau \geq t\}} \mathbb{E}\left[ X_t - X_{t-1} \mid \mathcal{F}_{t-1} \right]
$$

This equation completely resolves the proof for cases (b), (c), and (d):
\begin{itemize}
    \item \textbf{(b) Martingale:} $\mathbb{E}[X_t - X_{t-1} \mid \mathcal{F}_{t-1}] = 0$. Hence, $\mathbb{E}[X_{t \wedge \tau} \mid \mathcal{F}_{t-1}] = X_{(t-1) \wedge \tau}$. The stopped process is a martingale.
    \item \textbf{(c) Submartingale:} $\mathbb{E}[X_t - X_{t-1} \mid \mathcal{F}_{t-1}] \geq 0$. Because $\mathbf{1}_{\{\tau \geq t\}} \geq 0$, the second term is non-negative. Hence, $\mathbb{E}[X_{t \wedge \tau} \mid \mathcal{F}_{t-1}] \geq X_{(t-1) \wedge \tau}$. The stopped process is a submartingale.
    \item \textbf{(d) Supermartingale:} $\mathbb{E}[X_t - X_{t-1} \mid \mathcal{F}_{t-1}] \leq 0$. Thus, $\mathbb{E}[X_{t \wedge \tau} \mid \mathcal{F}_{t-1}] \leq X_{(t-1) \wedge \tau}$. The stopped process is a supermartingale.
\end{itemize}
\end{proof}

\chapter{Snell Envelope and Optimal Stopping}

The Snell envelope is a profound mathematical construct in stochastic analysis, directly formulated to solve optimal stopping problems. We have already seen its precursor intuitively in the dynamic programming equation of the American option.

\section{Definition and Construction}

Let $(\Omega, \mathcal{F}, \mathbb{P})$ be a generic probability space equipped with a discrete filtration $\mathbb{F} = (\mathcal{F}_{n})_{n=0, \dots, N}$. Let $(Z_n)_{n=0, \dots, N}$ be an $\mathbb{F}$-adapted process taking positive values. In our financial context, you should immediately think of $\mathbb{P}$ as the risk-neutral measure $\mathbb{Q}$, $n$ as the time step $t$, and $Z_n$ as the discounted payoff $\tilde{H}_t$, although the Snell envelope formulation is entirely general.

\begin{definition}[Snell Envelope]
The Snell envelope of the process $(Z_n)_{n=0, \dots, N}$ is the stochastic process $(E_n)_{n=0, \dots, N}$ defined recursively backwards in time by the following dynamic programming equations:
$$
\begin{cases}
E_N = Z_N, \\
E_n = \max\left( Z_n, \, \mathbb{E}\left[ E_{n+1} \mid \mathcal{F}_n \right] \right), \quad n = 0, \ldots, N-1.
\end{cases}
$$
\end{definition}

As derived intuitively in the first section regarding discounted prices, this construction naturally leads to a minimal bounding property.

\begin{theorem}[Smallest Dominating Supermartingale]
The Snell envelope $(E_n)_{n=0, \dots, N}$ is a supermartingale. Furthermore, it is the \textit{smallest} supermartingale that dominates $(Z_n)_{n=0, \dots, N}$ from above.
\end{theorem}

\begin{proof}
By definition of the Snell envelope, $E_n = \max( Z_n, \mathbb{E}[ E_{n+1} \mid \mathcal{F}_n ] )$. 
Since the maximum of two numbers is greater than or equal to both numbers, it trivially holds that:
\begin{enumerate}
    \item $E_n \geq \mathbb{E}[ E_{n+1} \mid \mathcal{F}_n ]$, proving $(E_n)$ is a supermartingale.
    \item $E_n \geq Z_n$, proving $(E_n)$ dominates the process $(Z_n)$ from above.
\end{enumerate}

To prove it is the \textbf{smallest} such process, let $(M_n)_{n=0,\dots,N}$ be \textit{any} other supermartingale that dominates $(Z_n)$ from above (i.e., $M_n \geq Z_n$ for all $n$). We proceed by backward induction down from $N$:
\begin{itemize}
    \item \textbf{Base Case ($n = N$):} $E_N = Z_N$. Since $M_N \geq Z_N$, we have $M_N \geq E_N$.
    \item \textbf{Inductive Step:} Assume $M_{n+1} \geq E_{n+1}$ almost surely. We must prove $M_n \geq E_n$.
    Because $(M_n)$ is a supermartingale, $M_n \geq \mathbb{E}[ M_{n+1} \mid \mathcal{F}_n ]$. 
    By the inductive hypothesis and the monotonicity of conditional expectations:
    $$ \mathbb{E}[ M_{n+1} \mid \mathcal{F}_n ] \geq \mathbb{E}[ E_{n+1} \mid \mathcal{F}_n ] $$
    So, $M_n \geq \mathbb{E}[ E_{n+1} \mid \mathcal{F}_n ]$. 
    By the dominating assumption, we also have $M_n \geq Z_n$. 
    Therefore, since $M_n$ is greater than both components inside the maximum:
    $$ M_n \geq \max\left( Z_n, \, \mathbb{E}\left[ E_{n+1} \mid \mathcal{F}_n \right] \right) = E_n $$
\end{itemize}
Thus, $M_n \geq E_n$ for all $n$, confirming $E_n$ is indeed the smallest dominating supermartingale.
\end{proof}

\section{The Optimal Stopping Problem}

The fundamental problem we are trying to solve is optimizing an expected payoff if we can stop at any valid stopping time. 

Let $\mathcal{V}(0, N)$ denote the collection of all possible stopping times taking values in $\{0, 1, \ldots, N\}$. We want to find a stopping time $\nu_0$ that maximizes the expected payoff evaluated at time $0$:
$$ \sup_{\nu \in \mathcal{V}(0,N)} \mathbb{E}\left[ Z_\nu \mid \mathcal{F}_0 \right] $$

This theoretical supremum represents the highest average value an optimal participant could extract from the game. How do we find this optimal stopping time $\nu_0$? The Snell envelope provides the exact recipe.

\begin{definition}[Optimal Stopping Time]
Define the random variable $\nu_0$ as the \textit{first} time the Snell envelope process equals the underlying payoff process:
$$ \nu_0 = \inf \{ n \geq 0 : E_n = Z_n \} $$
\end{definition}

\begin{lemma}
$\nu_0$ is a valid stopping time, and the stopped Snell envelope process $(E_{n \wedge \nu_0})_{n=0, \dots, N}$ is a true martingale.
\end{lemma}

\textbf{Intuition for the Martingale property:} Recall that $E_n = \max(Z_n, \mathbb{E}[E_{n+1}\mid \mathcal{F}_n])$. As long as we haven't stopped (i.e., for $n < \nu_0$), it means $E_n > Z_n$. Hence, the maximum evaluates strictly to the expectation:
$$ E_n = \mathbb{E}[E_{n+1} \mid \mathcal{F}_n] \quad \text{for } n < \nu_0 $$
This is the definition of a martingale! The Snell Envelope stops "decaying" (losing supermartingale value) while we hold out for the optimal time. 

\begin{theorem}[Resolution of the Optimal Stopping Problem]
The stopping time $\nu_0$ satisfies the optimal stopping problem. Evaluated at time 0, we have:
$$ E_0 = \mathbb{E}\left[ Z_{\nu_0} \mid \mathcal{F}_0 \right] = \sup_{\nu \in \mathcal{V}(0,N)} \mathbb{E}\left[ Z_\nu \mid \mathcal{F}_0 \right] $$
\end{theorem}

\begin{proof}
Let's break this down into two necessary steps:
\begin{enumerate}
    \item \textbf{Show that $E_0 \geq \mathbb{E}[Z_\nu \mid \mathcal{F}_0]$ for \textit{any} stopping time $\nu$:}
    Because $(E_n)$ is a supermartingale, Optional Stopping tells us that $\mathbb{E}[E_\nu \mid \mathcal{F}_0] \leq E_0$. 
    Furthermore, because the Snell envelope dominates $Z_n$ everywhere, $E_\nu \geq Z_\nu$. 
    Thus:
    $$ E_0 \geq \mathbb{E}[E_\nu \mid \mathcal{F}_0] \geq \mathbb{E}[Z_\nu \mid \mathcal{F}_0] $$
    Since this holds for any $\nu$, it must hold for the supremum:
    $$ E_0 \geq \sup_{\nu \in \mathcal{V}(0,N)} \mathbb{E}[ Z_\nu \mid \mathcal{F}_0 ] $$
    
    \item \textbf{Show that $\nu_0$ achieves this value exactly:}
    From the previous Lemma, $(E_{n \wedge \nu_0})$ is a martingale. Therefore, its expectation does not drift:
    $$ E_0 = \mathbb{E}[E_{\nu_0 \wedge N} \mid \mathcal{F}_0] $$
    But since $\nu_0$ can be at most $N$ (because $E_N = Z_N$, ensuring the set $\{n: E_n = Z_n\}$ is never empty), $\nu_0 \wedge N = \nu_0$.
    Thus:
    $$ E_0 = \mathbb{E}[E_{\nu_0} \mid \mathcal{F}_0] $$
    By definition of $\nu_0$, evaluated precisely at time $\nu_0$, $E_{\nu_0} = Z_{\nu_0}$. Therefore:
    $$ E_0 = \mathbb{E}[Z_{\nu_0} \mid \mathcal{F}_0] $$
\end{enumerate}
Combining both parts shows that $E_0$ is the supremum and $\nu_0$ is the optimal stopping time achieving that supremum.
\end{proof}

\subsection{Generalization to Intermediate Times}

This logic can be applied not just at time 0, but starting from any time $n$. We can define intermediate optimal stopping times $\nu_n$ indicating the optimal time to stop given we have already reached time $n$ without stopping.

We can define it as:
$$ \nu_n := \inf \left\{ k \geq n : E_k = Z_k \right\} $$

Analogous to the time 0 case, we find the dynamic value:
$$ E_n = \sup_{\nu \in \mathcal{V}(n,N)} \mathbb{E}\left[ Z_\nu \mid \mathcal{F}_n \right] = \mathbb{E}\left[ Z_{\nu_n} \mid \mathcal{F}_n \right] $$
where $\mathcal{V}(n, N)$ represents the collection of stopping times taking values in the future range $\{n, n+1, \ldots, N\}$.

This concludes the pure mathematical framework. The value of an optimally timed opportunity is exactly its Snell envelope. We will now apply this rigorously back to American Options.

\chapter{Pricing and Replication of American Options}

Having formally established the Snell envelope as the solution to the optimal stopping problem, we now return to our original financial objective: pricing and replicating American options in a discrete complete market.

\section{American Options Value Process}

The connection between the Snell envelope and American option pricing is immediate and direct. The dynamic programming principle introduced in the first section was simply the Snell envelope defined under the risk-neutral measure $\mathbb{Q}$.

\begin{proposition}[American Option Pricing]
Let $\mathcal{T}$ denote the collection of all valid stopping times taking values in $\{0, \dots, T\}$. 
The discounted value process $(\tilde{U}_t)_{t=0, \dots, T}$ of an American option defined by recursion:
$$
\begin{cases}
\tilde{U}_T = \tilde{H}_T \\
\tilde{U}_{t-1} = \max\left( \tilde{H}_{t-1}, \, \mathbb{E}^{\mathbb{Q}}\left[ \tilde{U}_t \mid \mathcal{F}_{t-1} \right] \right), \quad t = 1, \ldots, T
\end{cases}
$$
is exactly the Snell envelope of the discounted payoff process $(\tilde{H}_t)$.

Therefore, it flawlessly constitutes the solution to the optimal stopping problem for the option holder. The fair initial price of the American option, $U_0$, is given by:
$$
U_0 = \sup_{\tau \in \mathcal{T}} \mathbb{E}^{\mathbb{Q}}\left[ \frac{H_\tau}{S_\tau^0} \right]
$$
where $S_0^0 = 1$. The holder's optimal exercise policy is given by the stopping time $\nu_0 = \inf\{t \geq 0 : \tilde{U}_t = \tilde{H}_t\}$, or equivalently, the very first time the option's holding value drops exactly matching its immediate intrinsic exercise value.
\end{proposition}

\section{American vs. European Options}

A recurring question in quantitative finance is whether an American option should ever be exercised strictly before maturity, and how its value compares to an otherwise identical European option (which can only be exercised at maturity $T$).

Let us denote:
\begin{itemize}
    \item $V^a_t$: The value of an American option at time $t$ with payoff stream $(H_0, \ldots, H_T)$.
    \item $V^e_t$: The value of a European option at time $t$ with identical final payoff $H_T$.
\end{itemize}

\begin{proposition}[Value Comparison]
The following relationships hold between American and European options:
\begin{enumerate}[label=(\alph*)]
    \item \textbf{Early Exercise Premium:} The American option is never worth less than the European option at any time: 
    $$ V^a_t \geq V^e_t \quad \text{for all } t \in \{0, \ldotp \ldots, T\} $$
    This is intuitive: the American option gives the holder \textit{more} rights (early exercise), which must command a non-negative premium. Mathematically, for the European option, $T$ is just one specific stopping time within $\mathcal{T}$, hence the supremum over all $\mathcal{T}$ cannot be lower than the expectation evaluated at one specific stopping time $T$.
    \item \textbf{Condition for Equality:} If the current intrinsic payoff is always less than or equal to the European holding value, i.e., $H_t \leq V^e_t$ for all $t$, then early exercise is never optimal. Consequently, the options have exactly the same price process:
    $$ V^a_t = V^e_t \quad \text{for all } t $$
\end{enumerate}
\end{proposition}

\subsection{The Case of Call and Put Options}

This equality condition leads us to an extremely famous result regarding regular call options on non-dividend-paying stocks.

\begin{corollary}[Non-optimality of Early Exercise for Calls]
For the case of a standard American call option, the payoff is $H_u = (S_u - K)^+$. If the risk-free interest rate is non-negative ($r \geq 0$), then an American call option should never be exercised early. Hence, its value exactly matches the European call:
$$ V^a_t = V^e_t \quad \text{for all } t $$
\end{corollary}

\textbf{Intuition:} If you exercise it early, you receive $S_t - K$ immediately. However, if you instead sell the option, its time value (especially given interest $r \geq 0$ bounds the strike present value from below) ensures $V^e_t > S_t - K$. A rational investor would sell the option in the open market rather than exercise it and destroy that time value.

\begin{remark}
This result is \textbf{no longer true} for an American put option ($H_t = (K - S_t)^+$). Due to the time value of money, receiving $K$ today is worth strictly more than receiving $K$ at maturity. If the asset price drops to 0 immediately, the put holder wants their $K$ right now to start earning interest. Hence, an American put is strictly more valuable than a European put ($V^a_t > V^e_t$), confirming the existence of a positive early exercise premium. Similar deviations occur for calls on dividend-paying stocks.
\end{remark}

\section{Replication of American Options}

In a complete market, any European derivative can be replicated. Replicating American derivatives is far more complex due to the embedded optimal exercise choice. Let $U_t$ be the value of an American option at time $t$. 

\begin{proposition}[Seller and Buyer Replication Dynamics]
Assume the market is complete. Then the following properties describe replication:
\begin{enumerate}[label=(\roman*)]
    \item \textbf{Existence of a super-replicating strategy:} The Snell envelope's supermartingale property implies by the Doob-Meyer decomposition that there exists a trading strategy $(\xi^0_t, \xi_t)$ whose portfolio value process $V_t$ perfectly satisfies:
    $$ V_t = U_t + A_t $$
    where $A_t$ is a \textit{predictable, non-decreasing} consumption process, with $A_0 = 0$.
    
    \item \textbf{Seller's perspective (Perfect Replication):} From the seller's point of view, they must be prepared to pay out at any time the option buyer demands it. The seller can set up a portfolio costing $U_0$ initially and follow the strategy $(\xi^0_t, \xi_t)$. As long as the buyer holds the option, the seller's portfolio value $V_t \geq U_t$. If the buyer exercises optimally at $\tau$, the seller's portfolio covers the payoff perfectly ($V_\tau \geq U_\tau = H_\tau$).
    
    \item \textbf{Buyer's perspective (Optimal Exercise):} On behalf of the buyer, the optimal stopping time $\tau$ dictates they exercise exactly when the intrinsic payoff equals the option value: $H_\tau = U_\tau$. 
    Furthermore, they must not wait so long that the consumption process $A_t$ ticks upwards. They should exercise before or precisely when time decay completely erodes the option advantage over payoff:
    $$ \tau \leq \inf \{ k \geq 0 : A_{k+1} > 0 \} $$

    \item \textbf{Suboptimal Execution Penalty:} If the buyer does \textit{not} exercise at an optimal time (for instance, they hold it too long or exercise prematurely), the predictable process $A_t$ accumulates strictly positive value. Since the seller is executing a perfect super-replicating strategy, this discrepancy ($A_t > 0$) strictly becomes a risk-free monetary gain (arbitrage profit) for the seller at the expense of the irrational buyer.
\end{enumerate}
\end{proposition}

This completes our comprehensive mathematical review of American option pricing in discrete time, providing exact formulations based on backward induction, Snell envelopes, and optimal stopping algorithms.


\end{document}
